\documentclass{article}
\usepackage[T1]{fontenc}
\usepackage{lmodern}
\usepackage{graphicx}
\usepackage{amsmath,amssymb}
\usepackage[a4paper,margin=2cm]{geometry}
\usepackage[absolute,overlay]{textpos}
\setlength{\TPHorizModule}{1mm}
\setlength{\TPVertModule}{1mm}
\usepackage[indonesian]{babel}
\usepackage{hyperref}
\setlength{\emergencystretch}{2em}
% unit: Halaman 45

\begin{document}

% === Halaman 45 (python/native) ===

• Alice memilih bilangan acak k yang harus terletak di dalam selang [1, p-1]

• Cipherteks dari PM  adalah pasangan titik

• PC = [ (kB), (PM

 + kPB) ]


• Untuk mendekripsi, Bob mula-mula menghitung hasil kali titik pertama dari PC 

(yaitu kB) dengan kunci privatnya, b

• b  (kB)

• Bob kemudian mengurangkan titik kedua  dari PC dengan hasil kali di atas

• (PM + kPB) – [b(kB)] = PM + k(bB) – b(kB) = PM

• Bob kemudian men-decode PM untuk memperoleh pesan M

45

*) Sumber bahan: Debdeep Mukhopadhyay, Elliptic Curve Cryptography ,

 Dept of Computer Sc and Engg IIT Madras

\end{document}
